\documentclass[11pt,a4paper]{article}

\usepackage[T1]{fontenc}
\usepackage[utf8]{inputenc}
\usepackage[brazil]{babel}
\usepackage[margin=2.4cm]{geometry}
\usepackage{microtype}
\usepackage{booktabs}
\usepackage{caption}
\usepackage[hidelinks]{hyperref}

\setlength{\parindent}{0pt}
\setlength{\parskip}{5pt}
\renewcommand{\arraystretch}{1.12}
\newcommand{\arquivo}[1]{\texttt{#1}}

\begin{document}

\begin{center}
    {\small PUCRS \enspace $\cdot$ \enspace Sistemas Operacionais \enspace $\cdot$ \enspace Trabalho Prático 2}\\[6pt]
    {\LARGE\bfseries Experimentos de gerenciamento de memória}\\[8pt]
    Jayme Fortes, Rodrigo Gagliardi e Lucca Marques
\end{center}

\section{Arquivos de teste e método}

O experimento compara partições variáveis com Worst-Fit e Circular-Fit,
sistema buddy e paginação simples. Para observar o efeito de cada escolha,
fixamos a memória em 64 posições, começamos cada execução com a memória
vazia e usamos o modo contínuo. O simulador foi compilado com
\arquivo{javac --release 17}.

\begin{center}
\captionof{table}{Arquivos próprios que acompanham o código em \arquivo{tests/}.}
\label{tab:arquivos}
\small
\begin{tabular}{@{}lrl@{}}
\toprule
Arquivo & Linhas & Cenário \\
\midrule
\arquivo{teste1.txt} & 26 & Sequência com menos rejeições no Circular-Fit \\
\arquivo{teste2.txt} & 26 & Sequência com menos rejeições no Worst-Fit \\
\arquivo{teste3.txt} & 27 & Carga variada com alocações, liberações e acessos \\
\bottomrule
\end{tabular}
\end{center}

Os três arquivos incluem pedidos rejeitados por falta de espaço, repetição
de ID e liberação de processo inexistente. O \arquivo{teste3.txt} também
inclui um acesso fora do intervalo do processo na paginação e uma linha
malformada. Os exemplos distribuídos com o enunciado serviram para
conferir o comportamento dos modos antes dos experimentos.

\subsection*{Métricas}

\begin{itemize}
    \item A taxa de rejeição é o número de pedidos recusados por falta de
    espaço dividido pela soma de pedidos admitidos e rejeitados por falta
    de espaço. Erros de comando ficam fora desse denominador.
    \item A ocupação máxima mede o pico de posições reservadas. No buddy
    e na paginação, a sobra dentro de um bloco ou frame ocupado também
    conta como memória reservada.
    \item Nas partições variáveis, medimos a fragmentação externa como o
    total livre menos o tamanho do maior bloco livre. No buddy e na
    paginação, medimos a fragmentação interna como a soma das sobras
    dentro dos blocos ou frames reservados.
\end{itemize}

O simulador calcula as médias sobre o estado após cada comando reconhecido,
inclusive comandos que retornam erro ou aviso. Linhas vazias e malformadas
ficam fora da média. As medidas de fragmentação são dadas em posições.

\clearpage

\section{Worst-Fit e Circular-Fit}

Executamos os mesmos três arquivos nas duas políticas. O Worst-Fit busca
o maior bloco que comporte o pedido. O Circular-Fit retoma a busca a partir
do endereço seguinte à última alocação e usa o primeiro bloco adequado.

\begin{center}
\captionof{table}{Comparação das políticas de partições variáveis.}
\label{tab:politicas}
\small
\begin{tabular}{@{}llrrrrr@{}}
\toprule
Arquivo & Política & Admit. & Rej. & Taxa & Pico & Frag. ext. \\
\midrule
\arquivo{teste1} & Worst-Fit    &  7 & 7 & 50,0\% & 64/64 & 1,56 \\
\arquivo{teste1} & Circular-Fit & 11 & 3 & 21,4\% & 59/64 & 2,36 \\
\arquivo{teste2} & Worst-Fit    & 12 & 4 & 25,0\% & 60/64 & 4,68 \\
\arquivo{teste2} & Circular-Fit &  8 & 8 & 50,0\% & 60/64 & 2,32 \\
\arquivo{teste3} & Worst-Fit    &  9 & 3 & 25,0\% & 55/64 & 4,12 \\
\arquivo{teste3} & Circular-Fit &  9 & 3 & 25,0\% & 55/64 & 3,08 \\
\bottomrule
\end{tabular}
\end{center}

No \arquivo{teste1}, o Circular-Fit teve quatro rejeições a menos. No
\arquivo{teste2}, o Worst-Fit teve quatro a menos. Em ambos, a política
com menos rejeições teve maior fragmentação externa média. Essa média
resume os estados observados ao longo da execução, enquanto uma admissão
depende de haver, naquele instante, um bloco com tamanho suficiente para
o pedido. No \arquivo{teste3}, ambas admitiram nove pedidos, e o
Circular-Fit teve menor fragmentação externa média.

\section{Buddy e partições variáveis}

\begin{center}
\captionof{table}{Resultados do buddy sobre os mesmos arquivos.}
\label{tab:buddy}
\small
\begin{tabular}{@{}lrrrrr@{}}
\toprule
Arquivo & Admit. & Rej. & Taxa & Pico & Frag. int. \\
\midrule
\arquivo{teste1} &  9 & 5 & 35,7\% & 64/64 & 7,96 \\
\arquivo{teste2} & 10 & 6 & 37,5\% & 64/64 & 7,76 \\
\arquivo{teste3} &  7 & 5 & 41,7\% & 64/64 & 15,35 \\
\bottomrule
\end{tabular}
\end{center}

Nas partições variáveis, um pedido de 9 ocupa exatamente 9 posições.
Na liberação, a partição pode ser unida a qualquer vizinho livre contíguo.
No buddy, um pedido de 9 reserva um bloco de 16, com 7 posições de sobra.
A fusão no buddy tem um candidato determinado pela árvore e pode se
repetir até a raiz. Ela exige que os dois blocos sejam irmãos e estejam
livres, mesmo quando há outros blocos livres fisicamente adjacentes.

No \arquivo{teste3}, o buddy atingiu 64 posições reservadas e recusou cinco
pedidos. Cada política de partições variáveis recusou três. Os 15,35 do
buddy medem posições reservadas que ficaram sem uso dentro de blocos.
Os 4,12 e 3,08 das partições variáveis medem posições livres fora do maior
bloco livre. Cada indicador ajuda a explicar como o espaço fica
indisponível para determinados pedidos.

\clearpage

\section{Efeito do tamanho da página}

Executamos \arquivo{teste3.txt} com páginas de 4, 8 e 16 posições.
Para medir o tamanho das tabelas, contamos as entradas da tabela de
\arquivo{P2}, que pede 12 posições e é admitido nas três execuções.
Também registramos o maior número de entradas observado em uma tabela
individual, considerando todos os processos admitidos em cada execução.

\begin{center}
\captionof{table}{Paginação de \arquivo{teste3.txt} com memória de 64 posições.}
\label{tab:paginas}
\small
\begin{tabular}{@{}rrrrrrrr@{}}
\toprule
Página & Frames & Admit. & Rej. & Taxa & Frag. int. & P2 & Tabela máx. \\
\midrule
 4 & 16 & 8 & 4 & 33,3\% &  7,04 & 3 & 7 \\
 8 &  8 & 7 & 5 & 41,7\% & 13,27 & 2 & 4 \\
16 &  4 & 6 & 6 & 50,0\% & 26,27 & 1 & 1 \\
\bottomrule
\end{tabular}
\end{center}

Com páginas de 4 posições, \arquivo{P2} usa três páginas e ocupa
12 posições, sem sobra. Com páginas de 8 ou 16, reserva 16 posições,
com sobra de 4, e sua tabela passa a ter duas ou uma entrada.
Essa comparação mantém o tamanho do processo fixo. O máximo de cada
execução também depende de quais outros processos foram admitidos.

O pico foi de 64 posições nos três casos. Com páginas maiores,
mais posições ficaram reservadas nas últimas páginas dos processos
admitidos. Para esta carga, isso veio acompanhado de mais rejeições
e de tabelas menores. Entre os três tamanhos, 4 posições produziram
a menor taxa de rejeição. Uma página de 1 posição em uma rodada extra
admitiu nove pedidos e rejeitou três, com fragmentação interna zero.
Nesse caso, \arquivo{P2} precisou de 12 entradas e a maior tabela
observada teve 25, o que explicita o custo em entradas por processo.

\subsection*{Reprodução dos experimentos}

Execute a partir da raiz do projeto, com JDK 17 ou superior. No Windows,
\verb|.\start| compila e abre o simulador pelo PowerShell. No Linux e no
macOS, compile e inicie a aplicação com os comandos abaixo.

\begin{verbatim}
javac -encoding UTF-8 --release 17 -d out $(find src -name '*.java')
java -cp out memory.Main
\end{verbatim}

\textbf{Exemplo.} Para executar \arquivo{tests/teste1.txt} com Worst-Fit e
memória de 64 posições, passe os parâmetros ao iniciar o Java:

\begin{verbatim}
java -cp out memory.Main tests/teste1.txt --mode 1 --mem 64 \
    --policy worst --continuous
\end{verbatim}

No modo interativo, o mesmo teste usa as respostas \arquivo{1},
\arquivo{64}, \arquivo{1}, \arquivo{tests/teste1.txt} e \arquivo{n},
uma por vez. O resumo esperado é de 7 pedidos admitidos e 7 rejeitados
por falta de espaço. Para rodar os demais casos no modo interativo,
use as opções a seguir.

\begin{itemize}
    \item \textbf{Partições variáveis.} Digite \arquivo{1} para o modo,
    \arquivo{64} para a memória, \arquivo{1} para Worst-Fit ou \arquivo{2}
    para Circular-Fit, o arquivo (por exemplo, \arquivo{tests/teste1.txt})
    e \arquivo{n} para execução contínua. Repita com os três arquivos.
    \item \textbf{Buddy.} Digite \arquivo{2}, \arquivo{64}, o arquivo
    de teste e \arquivo{n}. Repita com os três arquivos.
    \item \textbf{Paginação.} Digite \arquivo{3}, \arquivo{64}, o tamanho
    da página (\arquivo{4}, \arquivo{8} ou \arquivo{16}),
    \arquivo{tests/teste3.txt} e \arquivo{n}. Repita para cada tamanho.
\end{itemize}

\clearpage

\section{Recomendações por tipo de carga}

Usamos memória de 64 posições nas três cargas adicionais da
Tabela~\ref{tab:cargas}. As sequências acompanham o código em
\arquivo{tests/}.

\begin{itemize}
    \item \arquivo{pequenos.txt} insere 20 processos de tamanho 3, libera
    \arquivo{S2}, \arquivo{S4}, \arquivo{S6} e \arquivo{S8} e insere quatro
    processos de tamanho 2.
    \item \arquivo{grandes-adjacentes.txt} e \arquivo{grandes-separados.txt}
    inserem \arquivo{L1} a \arquivo{L4}, todos de tamanho 16. O primeiro
    libera \arquivo{L2} e \arquivo{L3}. O segundo libera \arquivo{L1} e
    \arquivo{L3}. Ambos terminam com \arquivo{IN(L5,32)}.
\end{itemize}

Para reproduzir a tabela, siga os passos da seção anterior em modo
contínuo com cada um desses arquivos, usando o caminho \arquivo{tests/}
seguido do nome do arquivo. Execute as duas políticas do modo 1,
o modo 2 e o modo 3 com páginas de 1 e de 4 posições.

\begin{center}
\captionof{table}{Admissões e rejeições por falta de espaço em cada carga adicional.}
\label{tab:cargas}
\small
\begin{tabular}{@{}lccccc@{}}
\toprule
Carga & Worst-Fit & Circular-Fit & Buddy & Pág. 1 & Pág. 4 \\
\midrule
Muitos pequenos   & 24/0 & 24/0 & 20/4 & 24/0 & 20/4 \\
Grandes adjacentes &  5/0 &  5/0 &  4/1 &  5/0 &  5/0 \\
Grandes separados  &  4/1 &  4/1 &  4/1 &  5/0 &  5/0 \\
\bottomrule
\end{tabular}
\end{center}

\textbf{Muitos processos pequenos.} Recomendamos partições variáveis
quando há blocos contíguos suficientes. As duas políticas aceitaram todos
os 24 pedidos, com pico de 60 posições. O buddy e a paginação de 4
posições arredondaram cada pedido de 3 para 4 e recusaram quatro pedidos.
Se for necessário traduzir endereços, a paginação de 1 posição também
aceitou os 24, com uma entrada de tabela por posição.

\textbf{Poucos processos grandes.} Recomendamos paginação quando as
liberações tendem a deixar regiões separadas. Com páginas de 4, o pedido
de 32 coube nos dois cenários, sem fragmentação interna para esses
tamanhos. Com blocos adjacentes, ambas as políticas de partições variáveis
também atenderam ao pedido. O buddy recusou porque os dois blocos livres
de 16 pertenciam a pais diferentes na árvore e não podiam formar um
bloco buddy de 32.

\textbf{Tamanhos variados.} Para uma carga semelhante ao
\arquivo{teste3.txt}, recomendamos Circular-Fit quando a alocação precisa
ser contígua. Em relação ao Worst-Fit, manteve as três rejeições e reduziu
a fragmentação externa média de 4,12 para 3,08. Se a tradução de endereços
for necessária, páginas de 4 posições oferecem um ponto de partida para
equilibrar rejeições e tamanho das tabelas. Entre as páginas de 4, 8 e 16,
esse tamanho teve a menor taxa de rejeição. Os resultados opostos de
\arquivo{teste1} e \arquivo{teste2} mostram por que a escolha da política
deve ser validada com a carga esperada.

\subsection*{Uso de inteligência artificial}

O grupo usou Claude e Claude Code (Anthropic) no planejamento, na elaboração
de partes do código Java e na revisão. Ele também apoiou em ajustes na
interface de terminal, na validação do simulador e na redação deste relatório. As ferramentas foram usadas
conforme a seção 12 do enunciado do TP2.

\end{document}
